\section{Conflict Typing with a Decision Model}
\label{sec:method}

A typed decision model receives a \emph{state}, here the query and its numbered search results, and one or more
typed questions. For a \emph{choice} question it returns a probability for each option; for a \emph{yes/no}
question, the probability of \emph{yes}. Several questions over the same state are answered in one request. We
study two ways of asking for the conflict type.

\paragraph{Flat.} One choice question whose five options are the five categories. The instruction, the category
definitions and their examples are copied word for word from the prompt of \citet{cattan2025dragged}
(Appendix~\ref{app:prompts}); a test in our code fails if a single word differs. The predicted type is the
most probable option, and its probability is the model's confidence.

\paragraph{Tree.} The taxonomy is hierarchical, which suggests asking simpler yes/no questions in one request:
\emph{differ} (do the results give different answers?), \emph{compatible} (could all of them be correct?),
\emph{opinion} (is it a contested topic with opposing views?), \emph{outdated} (do they disagree because the answer
changed?), and \emph{misinformation} (are some results incorrect?). The five probabilities are combined in code into
a distribution over the categories. \textbf{Tree v1} follows the taxonomy's order: $P(\textit{No conflict}) =
1-d$, $P(\textit{Complementary}) = d\,c$, and the remaining mass $d(1-c)$ is split over the three causes in
proportion to their yes-probabilities, where $d$ and $c$ are the probabilities of \emph{differ} and
\emph{compatible}. \textbf{Tree v2}, chosen after inspecting v1 on the development split (\S\ref{sec:formulation}),
asks \emph{opinion} before \emph{compatible} and treats \emph{misinformation} as the factual conflict not explained
by time: $P(\textit{Opinions}) = d\,o$, $P(\textit{Complementary}) = d(1-o)c$, $P(\textit{Outdated}) =
d(1-o)(1-c)\,t$ and $P(\textit{Misinformation}) = d(1-o)(1-c)(1-t)$, with $o$ and $t$ the probabilities of
\emph{opinion} and \emph{outdated}. Both trees use the same five answers, so v2 costs nothing extra.

\paragraph{Cascade.} \jev answers every query. The fraction $f$ of queries with the lowest \jev confidence is sent
to an LLM, whose answer replaces \jev's. The threshold is a quantile of \jev's confidences and uses no labels, so the
whole accuracy--cost curve can be computed on the test split without leakage; we fixed $f=25\%$ and $f=50\%$ as the
reported operating points before computing any cascade.
